\section{An exact hierarchy with equal hypothesis priors}
\label{sec:equalprior91}
The strict separation does not require the prior in
\eqref{eq:gameprior90}.  Equal hypothesis priors lead to a different
optimal receiver readout and to another exact hierarchy.  The finite
weighted second-moment family is unchanged.

\subsection{A trace-norm bound}
Let $S$ be the swap on $\C^d\otimes\C^d$ and
$\Pi_\pm=(I\pm S)/2$.
\begin{lemma}[The centered filtered swap]\label{lem:centeredswap91}
For $d\ge2$, $A,B\succeq0$, and $L=\sqrt A\otimes\sqrt B$,
\begin{equation}\label{eq:centeredswap91}
 \|L(I-dS)L\|_1\le\left(d-\frac1d\right)\tr A\tr B.
\end{equation}
For nonzero factors equality holds precisely when $A$ and $B$ are
positive scalar multiples of $I$.  The same norm and bound hold for
rectangular factors having Gram matrices $A,B$.
\end{lemma}
\begin{proof}
Put $K=LSL$ and $f=\|\sqrt A\sqrt B\|_1$.
The spectrum calculation in Lemma~\ref{lem:swapfilter90} gives
$\|K\|_1=f^2$, including singular factors by continuity.
Since $I-dS=2\Pi_--(d-1)S$, the triangle inequality gives
\begin{align*}
 \|L(I-dS)L\|_1
 &\le 2\tr(L\Pi_-L)+(d-1)\|K\|_1\\
 &=\tr A\tr B-\tr(AB)+(d-1)f^2\\
 &\le\tr A\tr B+(d-1-d^{-1})f^2\\
 &\le(d-d^{-1})\tr A\tr B.
\end{align*}
The penultimate step is eigenvalue Cauchy--Schwarz,
$\tr(AB)\ge f^2/d$; the last is $f^2\le\tr A\tr B$.

Here is the equality argument, also supplying a detailed treatment of
singular nonzero factors in Lemma~\ref{lem:swapfilter90}.
Normalize $a=A/\tr A$, $b=B/\tr B$.  The variational formula for the
trace norm selects a unitary $U$ with
$\operatorname{Re}\tr(\sqrt a\sqrt b U)=\|\sqrt a\sqrt b\|_1$.
Both $\sqrt a$ and $\sqrt b U$ have Hilbert--Schmidt norm one.
Equality in Cauchy--Schwarz, with the phase chosen as above, forces
$\sqrt b U=\sqrt a$; multiplication by the adjoint gives $b=a$.
This implication uses no invertibility.  Equality in
$\tr(AB)\ge f^2/d$ then gives $\tr(a^2)=1/d$.  The $d$
nonnegative eigenvalues of $a$, including zeros, sum to one, so equality
in their Cauchy--Schwarz inequality forces all of them to equal $1/d$.
Thus nonzero singular factors cannot attain equality.  Since
$d-1-d^{-1}>0$, equality in \eqref{eq:centeredswap91} forces both
of these equalities.  Scalar factors attain the bound directly from
the eigenvalues $1-d$ on $\Pi_+$ and $1+d$ on $\Pi_-$.
Polar decompositions of rectangular factors preserve the nonzero
spectrum and add only zero eigenvalues, proving the last assertion.
\end{proof}

\subsection{Equal-prior optimization}
\begin{theorem}[Equal-prior operational hierarchy]\label{thm:equalprior91}
Let $(w_b,P_y^b)$ be any weighted exact second-moment family satisfying
\eqref{eq:basismoments90}.  Choose the scalar measurement $C_y=I/d$
with probability $1/2$, or choose the alternative family with probability
$1/2$ and then draw one basis index $b$ with law $w_b$.  In the latter
case the same device $E_y^b(t)=(1-t)I/d+tP_y^b$ is used at both calls.
For $d\ge2$ and $0\le t\le1$, the three classes of
Definition~\ref{def:classes91} have exact optimal success probabilities
\begin{equation}\label{eq:equalprior91}
 \begin{split}
 P_{\rm ca}^{\rm eq}&=\frac12+\frac{t^2(d-1)}{2d(d+1)},\\
 P_{\rm reset}^{\rm eq}&=\frac12+\frac{t^2(d-1)}{2d^2},\\
 P_{\rm all}^{\rm eq}&=\frac12+\frac{t^2}{2d}.
 \end{split}
\end{equation}
Both inequalities are strict for $t>0$, including the full-rank interval
$0<t<1$; all three values equal $1/2$ at $t=0$.
The reset optimum is attained with independent maximally entangled
input--reference pairs and no receiver feedback.  Its upper includes
all unit-reset receiver instruments and arbitrary receiver dimension.
The unrestricted optimum is attained by a normalized antisymmetric
two-input state and bounds every adaptive strategy.
If the alternative basis is instead redrawn independently at each
call, its two-slot averaged channel is the scalar product channel and
all three equal-prior values are $1/2$.
\end{theorem}
\begin{proof}
Put $a=t^2/[2d^2(d+1)]$.  Subtracting the alternative-weighted
second moment from the scalar-weighted product gives
\begin{equation}\label{eq:equalpayoff91}
 W_{yy}=a(I-dS),\qquad
 W_{yz}=-\frac{a}{d-1}(I-dS)\quad(y\ne z).
\end{equation}
Indeed the diagonal moment is
$(1-t^2)I/d^2+t^2(I+S)/[d(d+1)]$, and the off-diagonal
moment is $(1-t^2)I/d^2+t^2(dI-S)/[d(d^2-1)]$.
These substitutions prove \eqref{eq:equalpayoff91}.
Both $S$ and $\Pi_\pm$ are real in the chosen basis; their full
transpose does not alter these payoff blocks.  The gain over the
constant alternative decision is $\sum_{yz}\tr(T_{C,yz}W_{yz})$.

\emph{Classically adaptive upper.}
After a complete first readout, a branch has an input effect
$A\succeq0$, with $\sum_h A_{yh}=\rho_1$ for each $y$ and
$\tr\rho_1=1$.  The second single-call tester has normalization
$B\succeq0$, $\tr B=1$, and decision blocks $0\preceq T_z\preceq B$.
For $\tr A>0$ set $\widehat A=A/\tr A$.
Contracting \eqref{eq:equalpayoff91} with $A$ in the first slot gives
$a\tr A(I-d\widehat A)$ for $z=y$ and its negative divided by
$d-1$ for $z\ne y$.  For a Hermitian $Q$ and $0\preceq T\preceq B$,
$\tr(TQ)\le\tr(BQ_+)$.  Therefore the branch gain is at most
\[
 a\tr A\,\tr[B|I-d\widehat A|]\le a(d-1)\tr A.
\]
The final inequality follows from $0\preceq\widehat A\preceq I$ and
$d\ge2$.  The zero branch contributes zero.
Summing $\sum_{yh}\tr A_{yh}=d$ proves the first upper bound.
It covers arbitrary second-call references and decisions, not only
unassisted classical probes.  Norm closure preserves it.

To attain it, use the same fixed pure input $|u\rangle$ twice and
declare $C$ on unequal labels.  Under $C$ the equal-label probability is
$1/d$; under the averaged alternative it is
\[
 \sum_{by}w_b\langle u,E_y^b(t)u\rangle^2
 =\frac1d+\frac{t^2(d-1)}{d(d+1)}.
\]
This follows directly from the second moments.  The corresponding
Bayes gain is half their difference, namely $ad(d-1)$.

\emph{Unit-reset upper.}
Use Lemma~\ref{lem:resetnormal91} without any restriction on its
rectangular receiver spaces.  On a branch $(y,h)$ put
\[
 X=(C_{yh}\otimes D_{yh})(I-dS)(C_{yh}\otimes D_{yh})^*.
\]
The elementary identity
$\max_{0\preceq F\preceq I}\tr(FX)=\tr(X_+)$ and
\eqref{eq:equalpayoff91} give total branch gain at most
\[
 a\tr(X_+)+a\tr((-X)_+)=a\|X\|_1
 \le a(d-d^{-1})\tr A_{yh},
\]
since $\tr b_{yh}=1$.  There are $d-1$ opposite-sign output blocks,
which cancel their denominator.  Summing over $y,h$ gives
$a(d^2-1)=t^2(d-1)/(2d^2)$.
Countable outcomes and null histories use unnormalized positive tails.

For attainment, prepare two independent normalized maximally entangled
pairs $|\Omega\rangle/\sqrt d$.  On equal device labels declare $C$
when the references lie in $\Pi_-$; on unequal labels declare $C$
when the references lie in $\Pi_+$.  This is one final joint receiver
measurement, with complementary effect $I-F_C$ on each classical
block.  Its event probabilities are
\begin{equation}\label{eq:equalresetevents91}
 q_C=\frac{(d-1)(d+2)}{2d^2},\qquad
 q_b=\frac{(d-1)(d+2-2t^2)}{2d^2}
 \quad\hbox{for every }b.
\end{equation}
To verify them, the receiver block is
$(E_y^b(t))^{\mathsf T}\otimes(E_z^b(t))^{\mathsf T}/d^2$ and
$\tr[\Pi_\pm(A\otimes B)]=[\tr A\tr B\pm\tr(AB)]/2$.
Here $\tr E_y=1$,
$\tr E_y^2=t^2+(1-t^2)/d$, and
$\tr E_yE_z=(1-t^2)/d$ for $y\ne z$.
Counting the $d$ equal-label and $d(d-1)$ unequal-label blocks gives
\eqref{eq:equalresetevents91}.  Half the difference is the reset upper.

\emph{Unrestricted upper.}
By Proposition~\ref{prop:causalnormal91}, $T_{C,yz}\preceq\Xi_y$
and $\tr\Xi_y=1$.  The sum of the positive payoff parts for a fixed
first label is
\[
 (W_{yy})_++\sum_{z\ne y}(W_{yz})_+
 =a(d+1)\Pi_-+a(d-1)\Pi_+\preceq a(d+1)I.
\]
Thus the gain is at most $ad(d+1)=t^2/(2d)$.
Any normalized state supported on the nonzero antisymmetric input
subspace attains it by declaring $C$ on equal labels: their probability
is $1/d$ under $C$ and $(1-t^2)/d$ under every alternative.

The gaps are
\begin{equation}\label{eq:equalgaps91}
 P_{\rm reset}^{\rm eq}-P_{\rm ca}^{\rm eq}
 =\frac{t^2(d-1)}{2d^2(d+1)},\qquad
 P_{\rm all}^{\rm eq}-P_{\rm reset}^{\rm eq}
 =\frac{t^2}{2d^2}.
\end{equation}
Finally $\sum_bw_b\mathcal M_{E^b(t)}=\mathcal M_C$ as a channel
identity.  Two independent draws give its tensor square, hence the
same two-slot statistics even for adaptive testers.  This proves the
last assertion and the endpoint statements.
\end{proof}
The equal priors concern the two hypotheses, not a uniform prior on all
individual devices.  The result removes prior tuning, but still uses
a shared latent device.  It does not assert a gap for every fixed pair.
Together with Theorem~\ref{thm:resetvariational91}, it gives a general
optimization principle and an explicitly solved family, rather than
identifying these two levels of generality.

\subsection{Seven explicit qubit devices}
\label{sec:qubitexplicit91}
For $d=2$ write
\[
 X=\begin{pmatrix}0&1\\1&0\end{pmatrix},\quad
 Y=\begin{pmatrix}0&-i\\i&0\end{pmatrix},\quad
 Z=\begin{pmatrix}1&0\\0&-1\end{pmatrix}.
\]
The devices are $C=(I/2,I/2)$ and, for $A\in\{X,Y,Z\}$ and
$\epsilon\in\{-1,1\}$,
\[
 E^{A,\epsilon}(t)=
 \left(\frac{I+t\epsilon A}{2},\frac{I-t\epsilon A}{2}\right).
\]
The six alternatives have conditional weight $1/6$.
Both outcome orders are included: they cancel the first Pauli moments
and supply the required ordered second moments.  At $t=1$ they are the
six ordered Pauli eigenbasis measurements.  For the prior of
Theorem~\ref{thm:operationalhierarchy90}, $C$ has mass $2/5$ and each
alternative mass $1/10$; for Theorem~\ref{thm:equalprior91}, these masses
are $1/2$ and $1/12$.

For reset acquisition prepare
\[
 |\phi^+\rangle_{I_1R_1}\otimes|\phi^+\rangle_{I_2R_2},\qquad
 |\phi^+\rangle=(|00\rangle+|11\rangle)/\sqrt2.
\]
The inputs are independent.  The two retained references are measured
jointly only after both calls.  Let
$|\psi^-\rangle=(|01\rangle-|10\rangle)/\sqrt2$,
$\Pi_-=|\psi^-\rangle\langle\psi^-|$ and $\Pi_+=I_4-\Pi_-$
on $R_1\otimes R_2$.
The receiver block for any two devices is $E_y^{\mathsf T}\otimes
F_z^{\mathsf T}/4$.  Under $J^\sharp$, the decision tester block is
$F_{C,yz}^{\mathsf T}/4$, where
\[
 F_{C,yz}=\mathbf1_{y=z}\Pi_-
 \quad\hbox{for the prior }2/5,
 \qquad
 F_{C,yz}=\begin{cases}\Pi_-,&y=z,\\\Pi_+,&y\ne z\end{cases}
 \quad\hbox{for equal priors}.
\]
The other block is $(I_4-F_{C,yz})^{\mathsf T}/4\succeq0$.
Their sum is $I_4/4$ for each $(y,z)$, with partial trace
$I_2/2$ over $I_2$ and trace one.  This verifies the causal normalization
in the order $I_1,Y_1,I_2,Y_2$, using the canonical permutation to group
the input copies.  The transpose is required for arbitrary complex
receiver effects, even though the displayed antisymmetric projectors
are real.

At $t=1$, the first reset event has probabilities $1/8$ under $C$ and
zero under every alternative.  Its Bayes score is
$(2/5)(1/8)+(3/5)=13/20$.
The equal-prior reset event has probabilities $1/2$ and $1/4$,
respectively, giving $\frac12(\frac12)+\frac12(\frac34)=5/8$.
For unrestricted acquisition put the normalized state $|\psi^-\rangle$
on $I_1\otimes I_2$, not on the references, and declare $C$ on equal
labels.  The probabilities are $1/2$ and zero, yielding $4/5$ for the
first prior and $3/4$ for equal priors.  For classical adaptation,
the biased-prior optimum is the constant decision $3/5$.
For equal priors, two inputs $|0\rangle$ and the unequal-label event
have probabilities $1/2$ under $C$ and $1/3$ averaged over alternatives,
giving $7/12$.  Consequently the two exact triples on the same devices are
\[
 \frac35<\frac{13}{20}<\frac45,
 \qquad
 \frac7{12}<\frac58<\frac34.
\]
The class upper bounds, not these evaluations alone, prove optimality.

\subsection{An open neighborhood of the equal-prior experiment}
\begin{corollary}[Channel, prior, and ensemble stability]\label{cor:gamestability91}
Fix the family and $t>0$.  Replace the scalar device and each entire
basis measurement channel individually by channels within unhalved
diamond distance $\delta$.  The perturbations may depend on the basis
index and may change all its effects, but must form legal channels on
the same interfaces.  Keep the once-drawn-index rule.  Replace the
scalar prior $1/2$ by $q$ and the conditional basis law $w$ by $\nu$,
and put
\[
 \eta=\tfrac12\sum_b|\nu_b-w_b|,\qquad
 \kappa=\delta+|q-\tfrac12|+\eta.
\]
For each of the three fixed operational classes,
$|\widetilde P_{\rm class}-P_{\rm class}^{\rm eq}|\le\kappa$.
Thus both class gaps remain strict when
$2\kappa<t^2(d-1)/[2d^2(d+1)]$.
If the exhibited reset protocol has, in addition, a uniform
reference-complete diamond certificate of preparation defect whose
hard pathwise sum is at most $\varepsilon$, its implemented score
still exceeds the perturbed classical optimum provided
\[
 2\kappa+\varepsilon/2<\frac{t^2(d-1)}{2d^2(d+1)}.
\]
\end{corollary}
\begin{proof}
For a fixed protocol and device index, the two-call hybrid bounds
output trace distance by $2\delta$.  Normalization makes the change
of any decision probability at most $\delta$, with the factor $1/2$
for unhalved trace distance.  Changing the binary prior changes success
by at most $|q-1/2|$; changing the alternative law changes an average
of $[0,1]$-valued successes by at most $\eta$.  The bound is uniform
in the protocol, so taking suprema preserves it in both directions.
Subtract and use \eqref{eq:equalgaps91}.
The preparation hybrid changes each implemented output by at most
$\varepsilon$, hence its success by at most $\varepsilon/2$.
This last statement concerns a certified implementation of the displayed
reset protocol, not an enlargement of the ideal optimization class.
\end{proof}

\subsection{Quantitative rigidity of near-optimal reset strategies}
The equality condition for the centered swap has a stable form.  It
controls the input Gram matrices in a chosen purified normal form,
rather than specifying unique instruments, dilations, or receiver
readouts.  For $d\ge2$ set
\[
 c_d=d-1-d^{-1}>0,\qquad
 K_d=d\left(1+\frac{9}{2c_d}\right).
\]
\begin{lemma}[Stable centered-swap equality]\label{lem:stableswap91}
For density matrices $a,b$ on $\C^d$, let
\[
 \Delta(a,b)=d-d^{-1}
  -\| (\sqrt a\otimes\sqrt b)(I-dS)
                         (\sqrt a\otimes\sqrt b)\|_1.
\]
Then $\Delta(a,b)\ge0$ and
\begin{equation}\label{eq:stableswap91}
 \|a-I/d\|_1^2\le K_d\Delta(a,b),\qquad
 \|b-I/d\|_1^2\le K_d\Delta(a,b).
\end{equation}
No invertibility is required.
\end{lemma}
\begin{proof}
Put $f=\|\sqrt a\sqrt b\|_1$, $\eta=1-f^2$, and
$e=\tr(ab)-f^2/d$.  The proof of Lemma~\ref{lem:centeredswap91}
gives $\eta,e\ge0$ and
\begin{equation}\label{eq:swapdefectdecomposition91}
 \Delta(a,b)\ge c_d\eta+e.
\end{equation}
Choose a unitary $V$ extending the right polar factor so that
$Z=\sqrt a\sqrt b V=\sqrt{\sqrt a b\sqrt a}\succeq0$.
This is possible also when the matrices are singular.  The
Hilbert--Schmidt norm, denoted by $\|\cdot\|_2$, satisfies
\[
 \|\sqrt b V-\sqrt a\|_2^2=2(1-f),\quad
 \|Z-a\|_2\le\sqrt{2\eta},\quad
 \|Z-fI/d\|_2^2=e.
\]
Here $\|\sqrt a\|_{\rm op}\le1$ and $1-f\le\eta$.
Since also $1-f\le\sqrt\eta$, the triangle inequality yields
\[
 \|a-I/d\|_2\le(\sqrt2+d^{-1/2})\sqrt\eta+\sqrt e.
\]
Weighted Cauchy--Schwarz bounds its square by
$[1+(\sqrt2+d^{-1/2})^2/c_d](c_d\eta+e)$.
For $d\ge2$, $(\sqrt2+d^{-1/2})^2\le9/2$; and
$\|T\|_1^2\le d\|T\|_2^2$ for a $d\times d$ matrix $T$.
These facts and \eqref{eq:swapdefectdecomposition91} prove the first
estimate.  Interchanging $a,b$ proves the second, since swap
conjugation leaves the centered-swap norm unchanged.
\end{proof}

\begin{corollary}[Near-optimal acquisition marginals]\label{cor:resetstability91}
Fix the equal-prior experiment of Theorem~\ref{thm:equalprior91}
with $t>0$.  Suppose a unit-reset strategy achieves reward
$v\ge P_{\rm reset}^{\rm eq}-\varepsilon$.
Choose its purified normal form from Lemma~\ref{lem:resetnormal91}.
On nonzero branches put
\[
 a_{yh}=A_{yh}/\tr A_{yh},\qquad
 \mu_{yh}=\tr A_{yh}/d.
\]
Then $\sum_{yh}\mu_{yh}=1$ and
\begin{equation}\label{eq:resetstability91}
 \sum_{y,h}\mu_{yh}
 \left(\|a_{yh}-I/d\|_1^2+\|b_{yh}-I/d\|_1^2\right)
 \le \frac{4K_d d(d+1)}{t^2}\,\varepsilon.
\end{equation}
For a public mixture the same conclusion holds after adjoining the
seed and multiplying these weights by its law.  Finite and countable
recorded instruments are both allowed.
\end{corollary}
\begin{proof}
Write $\alpha=t^2/[2d^2(d+1)]$ for the coefficient in
\eqref{eq:equalpayoff91}.  The complete reset upper, before applying
the centered-swap bound, gives
\[
 v\le\tfrac12+\alpha\sum_{yh}\tr A_{yh}
          [d-d^{-1}-\Delta(a_{yh},b_{yh})].
\]
Purification does not obstruct recovering the original decision rule,
so this inequality bounds its achieved reward as well.  The common
Gram identity gives $\sum_{yh}\tr A_{yh}=d$.  Subtract from the
attained optimum to obtain
$\sum_{yh}\mu_{yh}\Delta(a_{yh},b_{yh})
\le\varepsilon/(\alpha d)$.
Apply Lemma~\ref{lem:stableswap91} to both factors and substitute
$\alpha$.  Positive trace tails justify countable sums.  For a public
mixture, apply the same upper to each component and average before
subtracting its reward from the common optimum.
\end{proof}
In particular, the total $\mu$-weight on which either normalized Gram
matrix is at trace distance at least $u>0$ from $I/d$ is at most
$4K_d d(d+1)\varepsilon/(t^2u^2)$.  These are Gram weights, not
probabilities under a selected hypothesis.  At zero loss the result
recovers branchwise isotropy.  The estimate is nontrivial at the
relative scale $\varepsilon/t^2$; it does not turn fixed calibration
error into vanishing-signal robustness.  It asserts neither uniqueness
of a dilation nor a device-independent certification of the reset cut.
